%% ma: L11-B15-0001
%% lop: 11
%% chuong: 5
%% bai: 15
%% dang: 11.15.1
%% loai: TL
%% muc: TH
%% dap_an: "-"
%% nguon: "tai-lieu-hoc-tap-mon-toan-lop-11-theo-chuong-trinh-giao-duc-pho-thong-2018.pdf (Lớp Toán Thầy Duy, Phạm Lê Duy, THPT Chu Văn An), Chương 3 Bài 1, Dạng 1, Ví dụ 1 (tr.306)"
%% kien_thuc: "Lớp 11 Bài 15 (giới hạn của dãy số)"
%% trang_thai: chinh-thuc
%% ghi_chu: "Quy ước kho: chứng minh bằng định nghĩa lim(u_n − L) = 0 và các định lí giới hạn (SGK Bài 15), không dùng ngôn ngữ epsilon–N. Đã duyệt 03/11/2026."
\begin{ex}
Chứng minh rằng

a) $\lim\left(\dfrac{-n^3}{n^3+1}\right)=-1$;

b) $\lim\left(\dfrac{n^2+3n+2}{2n^2+n}\right)=\dfrac12$.
\loigiai{\sloppy a) $\dfrac{-n^3}{n^3+1}-(-1)=\dfrac{1}{n^3+1}\to0$, nên giới hạn bằng $-1$.

b) \[\dfrac{n^2+3n+2}{2n^2+n}-\dfrac12=\dfrac{5n+4}{2(2n^2+n)}=\dfrac{\frac5n+\frac4{n^2}}{2\left(2+\frac1n\right)}\to0,\] nên giới hạn bằng $\dfrac12$.}
\end{ex}
